\chapter{突变：错误、材料还是创新}

\begin{kaichang}
秋天去果园，你可能会撞见一件怪事。同一棵苹果树，施的是同样的肥，晒的是同样的太阳，绝大部分枝条结的果子青里透红，唯独有一根枝条，上面每一个苹果都是深红色的，红得发亮，而且个头更大、成熟更早。不是虫子咬的，不是鸟啄的，也不是谁偷偷给它开了小灶——去年它是这样，今年还是这样，砍一枝下来嫁接别的树上，结出来的果子照样深红。

先猜一猜：这根枝条发生了什么？是土壤里多了什么东西？是这棵树的“心情”变了？还是它身体里的某份“文件”被改动了？

提示一句：我们今天吃的很多水果，包括好几种著名的红苹果，祖先正是果园里这样一根“不合群”的枝条。这一章结束时，我们会回来把这件事彻底讲清楚。
\end{kaichang}

\section{满世界都是“不一样的副本”}

先把眼光从苹果树移开，你会发现这类事到处都是：

\begin{itemize}[itemsep=2pt]
\item 花市里卖的“金边吊兰”“花叶绿萝”，叶子上带着白色或黄色的条纹和斑块，同一盆花里，有的叶子全绿，有的叶子半绿半白，偶尔还冒出一整片纯白的叶子。
\item 动物养殖场里偶尔会生出白色的个体：白化的兔子、白化的小鼠，甚至白化的鹿。它们的父母毛色完全正常。
\item 医院里，某种抗生素开头能把一批细菌杀得干干净净，用过几年之后，同样的药、同样的剂量，却治不动同样的感染了。
\item 我们每个人自己也是一例：你大约有三百万个位点上的 DNA 字母和你的同班同学不一样——当然，这些差异绝大多数来自父母，但其中还有几十个，是你父母双方的 DNA 里都找不到、只在你身上第一次出现的。
\end{itemize}

这些现象看上去毫不相干：水果、叶子、细菌、你自己。但把它们摆在一起，会冒出同一个问题——\textbf{复制为什么不是完美的}？第 67 章讲过，细胞分裂前要把整套 DNA 照抄一遍，几十亿个字母，抄得又快又准。可“又快又准”不等于“一字不差”。这一章的主角，就是抄本上和原件对不上的那些地方。它有一个你可能会觉得吓人的名字：\textbf{突变}。

请先记住一件事：在生物学里，“突变”是一个中性词，意思就是 DNA 序列发生了可遗传的改变。它既不天然等于“错误”，也不天然等于“进步”。它是错误、是材料、还是创新，要看场合——这正是本章标题的问题，我们会一路问到底。

\section{镜头推进：DNA 字母被怎样改动}

把镜头推回分子层。第 66 章说过，DNA 的全部信息写在四种碱基——腺嘌呤（A）、胸腺嘧啶（T）、胞嘧啶（C）、鸟嘌呤（G）——排成的长链里；第 67 章说过，复制就是照着旧链配出新链。所谓突变，就是这串字母的排布被改动了。改动的方式有好几种，尺度差别很大，先分成两类。

第一类是\textbf{小尺度改动}，只动几个字母：

\begin{itemize}[itemsep=2pt]
\item \textbf{替换}：一个字母被换成另一个，比如某个位置上的 A 变成了 G。好比抄写时写了个错别字。
\item \textbf{插入}：多抄了一个或几个字母。好比抄写时手一抖，重复写了某个字。
\item \textbf{缺失}：少抄了一个或几个字母。好比抄写时漏了字。
\end{itemize}

第二类是\textbf{大尺度改动}，动的是一整段文字，涉及几千到几百万个字母：

\begin{itemize}[itemsep=2pt]
\item \textbf{重复}：一整段序列被多抄了一份，染色体上同一篇“文章”出现了两次甚至更多次。
\item \textbf{倒位}：一整段序列断了下来，翻个面又接了回去，字母一个不少，顺序整个颠倒。
\item \textbf{易位}：一整段序列从一条染色体上断下来，接到了另一条染色体上，像把一页纸装订进了另一本书。
\end{itemize}

为什么要分得这么细？因为改动的方式决定了后果的性质。一个错别字可能无伤大雅，也可能改变整句话的意思；而一页纸装订错了书，影响的可能是好几章的内容。要判断后果，就得回到这些字母的“读法”上去——第 69 章讲过的三联体密码。

\section{一字之差，命运各异}

回忆一下第 69 章的规则：蛋白质合成时，机器沿着信使 RNA 三个字母一组地读，每组是一个密码子，对应一种氨基酸（或“停止”信号）。全文 64 个密码子对应 20 种氨基酸，所以有好几个密码子说的是同一种氨基酸——密码表是“简并”的，像好几个不同的代号指向同一个人。

有了这条规则，同一个字母的改动就能分出四种完全不同的命运。拿一段真实序列来演示——下面这一段来自人体血红蛋白的 \(\beta\) 链基因的开头附近：

\begin{itemize}[itemsep=2pt]
\item \textbf{沉默突变}：某个密码子的第三个字母变了，但新旧密码子对应同一种氨基酸。蛋白质一个氨基酸都没变，一切照旧。就像把“茴香豆”写成“茴香荳”，读音意思全没变。这是最常见的结果之一。
\item \textbf{错义突变}：改动后的密码子换了另一种氨基酸。蛋白质链条上有一个零件被换了。后果可大可小：换在不碍事的位置，可能毫无影响；换在关键位置，整台机器散架。
\item \textbf{无义突变}：改动后的密码子变成了“停止”信号。蛋白质合成到一半就提前收工，多半造出一截没用的半成品。
\item \textbf{移码突变}：插入或缺失的字母数不是 3 的倍数。这下不是改错一个词的问题了——阅读框从出错位置开始整体错位，后面每一组三个字母都被重新划分，一整段文字全军覆没。好比一句话里漏了一个字，后面每个词都读岔了。
\end{itemize}

现在你能明白为什么大多数突变“没什么大不了”了。第一层缓冲是密码表的简并：不少替换是沉默的。第二层缓冲是基因组的结构：人类 DNA 里，直接编码蛋白质的序列只占约百分之一多一点，大量突变落在不编码的区域，很多位置并不造成看得出的差别。第三层缓冲是蛋白质本身的宽容度：换上一个性格相近的氨基酸，机器常常照转不误。

但缓冲不是万能的。剩下的那一小部分突变，刚好砸在要害上——而“要害”是什么样子，下一节我们用一个真实的序列来看。

\section{把突变写下来：一段 DNA 的几种命运}

符号现在登场。血红蛋白 \(\beta\) 链的第 6 个氨基酸，正常密码子（写在 DNA 编码链上）是 GAG，对应谷氨酸——一个带负电、亲水的氨基酸。围绕这一个密码子，我们可以把上一节的四种命运各演一遍：

\begin{table}[htbp]
\centering
\begin{tabular}{llll}
\toprule
突变类型 & 序列变化 & 蛋白质后果 & 实际影响 \\
\midrule
正常 & GAG & 第 6 位是谷氨酸 & 血红蛋白正常 \\
沉默 & GAG 变 GAA & 还是谷氨酸 & 无影响 \\
错义 & GAG 变 GTG & 变成缬氨酸 & 镰刀型贫血 \\
无义 & GAG 变 TAG & 提前终止 & 链被截短，失能 \\
移码 & 前面插入 1 个字母 & 从此全部读岔 & 基本失能 \\
\bottomrule
\end{tabular}
\end{table}

表里的错义突变 GAG 变 GTG 是真实存在的：它把亲水的谷氨酸换成了疏水的缬氨酸，让血红蛋白分子表面多了一块“怕水的补丁”。脱氧时这些分子会黏在一起拉成长纤维，把红细胞撑成镰刀形。只改了一个字母，后果是一种伴随人类几千年的遗传病。这是“突变是错误”最有力的证据之一。

可故事还有另一半。只从父母一方继承这个突变（另一个 \(\beta\) 基因正常）的人，平时几乎健康，而在疟疾肆虐的地区，他们的红细胞反而让疟原虫难以生存——疟疾是那里最大的儿童杀手之一。于是同一个突变，在纯合时是重病，在杂合时是护身符，在没有疟疾的地区又变回纯粹的负担。\textbf{同一个字母，是错误还是创新，由环境投票决定}。这张票怎么投、投完怎样改变一个群体的基因组成，是第 76 章的主题，这里先把现象记住。

\section{错误从哪来，谁在校对}

突变有两个来源。第一个来源防不胜防：\textbf{复制本身就会出错}。DNA 聚合酶每秒钟能装配几百个字母，再精密的机器也有手滑的时候；碱基分子自己也会偶尔“变身”（第 66 章提过，碱基存在稀有的互变异构形态，配对错认），化学结构上的微小摆动就能骗过配对规则。这类错误叫\textbf{自发突变}——不需要任何外部原因，物理和化学规律本身就保证它会发生。

第二个来源是\textbf{诱变因素}，即从外面加大出错概率的东西：

\begin{itemize}[itemsep=2pt]
\item \textbf{紫外线}：阳光中的紫外线能让同一条链上相邻的两个胸腺嘧啶（T）黏成一个“疙瘩”（二聚体），复制机器读到这里就卡壳、乱猜。晒伤的代价不只是脱皮，而是皮肤细胞里实实在在增加的一批突变。防晒是货真价实的“防突变”，而且和剂量有关——晒多久、多强的太阳，决定风险的大小。
\item \textbf{某些化学物质}：烟草烟雾里的多环芳烃、霉变花生里的黄曲霉毒素、腌菜烤肉里可能形成的亚硝胺，都能直接损伤碱基或干扰复制。这里同样永远是剂量问题：偶尔吃一次烧烤和每天抽一包烟，风险差着几个数量级。
\item \textbf{电离辐射}：X 光、放射性物质放出的高能粒子能直接打断 DNA 链，断口重接时就可能发生缺失、倒位、易位这类大尺度改动。
\end{itemize}

好在细胞不是裸奔的。第 67 章讲复制时埋过一个伏笔：DNA 聚合酶一边抄一边\textbf{校对}，发现刚安上去的字母不配，就退回去拆掉重安，这一步就把错误率压低了约一百倍。复制结束后还有第二道防线：\textbf{错配修复}系统像质检员一样沿着新链巡查，揪出漏网的错配。此外还有针对紫外线疙瘩、化学损伤的\textbf{切除修复}——把受伤的一小段切掉，照另一条链重新合成补上。一层层防线叠下来，最终净错误率低到大约每抄十亿个字母才漏掉一个。

但“低”不是“零”。而且修复系统自己也是蛋白质，也是基因编码的——编码修复系统的基因同样可能突变。有一种罕见的遗传病叫着色性干皮病，患者的切除修复系统天生失灵，普通人晒太阳能修好的损伤他们修不好，于是对阳光极度敏感，皮肤癌风险极高。这个病本身就是“修复系统真的在工作”的活证据。

\begin{shiyan}
\textbf{用一次抄写模拟突变与校对}（安全，只需纸笔）

材料：一段 100 字左右的课文、纸、笔、计时器、一位同学。

步骤：① 你在一分钟内尽量快地抄写这段课文；② 和原文逐字比对，数出三类错误：写错的字（替换）、漏掉的字（缺失）、多写的字（插入）；③ 把这篇抄件交给同学，请他再抄一遍，同样统计；④ 最后两人互换“校对员”角色，把对方第二遍的抄件逐字订正一遍，再数还剩下几个错。

观察点：你的“突变率”大约是每抄多少个字出一个错？三类错误各占多少？经过“校对修复”之后，剩余错误率下降了多少倍？如果中间漏掉一个字，后面的字还读得通吗——这就是移码为什么格外致命。

安全提示：纯纸笔活动，无风险。唯一的要求是诚实——别偷偷改数据，科学家的名声就毁在这种事上。
\end{shiyan}

\begin{qiaoliang}
\textbf{你身上带着多少个“全新”的突变？} 来算一笔真实的账。人类单倍体基因组约 $3\times 10^{9}$ 个字母，经过校对加修复之后，每代每位的净突变率约 $1\times 10^{-8}$ 到 $2\times 10^{-8}$（这是直接比较父母和孩子的基因组测出来的，不是猜的）。乘一乘：
\[
3\times 10^{9}\ \times\ 1.5\times 10^{-8}\ \approx\ 45
\]
也就是说，你父母各贡献的一套基因组里，加起来大约有几十到一百个位点（常见估算是 60 到 100 个）上的字母，是在他们身上都没有、在你这里第一次出现的。全世界 80 亿人，每人几十个新突变——仅仅人类这一代，基因组里几乎每个能变的位点都被“试”过了不止一遍。再想想你身体里此刻约 $3\times 10^{13}$ 个细胞，每个细胞分裂时都要重抄一遍——\textbf{你本人就是一座不停产生新序列的图书馆}，只不过其中绝大多数“新书”无人问津，也不造成伤害。
\end{qiaoliang}

\section{科学家怎样知道突变不是“按需生产”的}

读到这里，一个很自然的想法可能在你脑子里冒头：细菌被抗生素杀死了那么多，活下来的产生了耐药性——这会不会是细菌“感受到”药物之后，有针对性地把自己改造了？需要在前，突变在后？如果是这样，突变就是“创新”最直接的证据：生命想要什么，就变出什么。

这个想法太重要了，必须接受实验的审判。

\begin{zhengju}
1943 年，遗传学家卢里亚和德尔布吕克设计了一个漂亮的实验，后来被称为\textbf{波动实验}。当时争论的问题是：细菌对噬菌体（一种吃细菌的病毒）的抗性，到底是接触病毒后才被“诱导”出来的，还是在接触之前就已经随机产生了？

两种假说给出的预言不一样。如果抗性是接触病毒后才诱导的，那么每个细菌“中招变身”的机会均等，取很多份细菌去测，抗性是后来才有的，各份里的抗性细菌数目应该像均匀撒盐：每份差不多，数量围绕一个平均值小幅波动。如果抗性是接触前就随机突变出来的，那么突变发生在培养过程的哪个时刻就很关键——第一代就突变的那支细菌，到实验结束时已经把抗性传给了成千上万个后代；最后一刻才突变的，只留下一两个后代。各份之间的差异应该极其悬殊，有的份一个抗性菌都没有，有的份成千上万。

实验做法很简单：把同一批细菌分成几十个小管，各自独立培养一段时间，然后分别涂到铺满噬菌体的培养皿上，数每个小管长出多少个抗性菌落。结果一边倒：小管与小管之间的菌落数差异巨大，波动远远超出“均匀撒盐”允许的范围。结论无可回避：\textbf{抗性突变在接触噬菌体之前就已经随机发生了，病毒只是把已有的抗性个体筛了出来}。后来人们把细菌的 DNA 直接测序，在分子层面一次又一次证实了同样的图景：突变先于选择，不朝向任何目标。

这个实验值得你记住的不只是结论，还有方法：两位科学家没有争论“生命有没有目的”这种哲学问题，而是把它翻译成了一个可以数数的统计问题——波动大还是波动小。哲学分歧一旦变成数字分歧，答案就自己走出来了。

顺便补一段前史：早在 1927 年，马勒就用 X 射线照射果蝇，让突变率提高了一百多倍，证明突变的速率可以被外界因素改变——这是诱变因素的第一次实锤，也为后来的辐射育种埋下了种子。
\end{zhengju}

\section{这个模型的边界在哪里}

“突变是 DNA 抄写和损伤留下的改动”这个模型很好用，但有几条边界你要知道。

\textbf{第一，突变不是均匀撒布的。} 基因组里有些位置格外容易出错（比如某些碱基组合是“热点”），有些区域被修复系统照顾得更周到，不同物种、不同性别的突变率也不一样。“随机”指的是不朝向需要，而不是处处等概率。

\textbf{第二，突变发生在哪个细胞里，命运完全不同。} 长在你身体普通细胞里的突变叫\textbf{体细胞突变}：它只影响那个细胞及其后代，不会传给你的孩子。开篇那根红苹果枝条就是体细胞突变——所以它只能靠嫁接、扦插这类无性繁殖保留下来。体细胞突变累积多了也有代价：癌症的本质，就是一个体细胞里若干关键基因接连被突变击中，挣脱了分裂的刹车。而发生在生殖细胞（精子、卵子）里的\textbf{生殖系突变}，会进入下一代的每一个细胞，成为可遗传的变异——进化真正关心的原材料，主要是这一类。

\textbf{第三，“有害还是有利”没有绝对答案。} 绝大多数改变蛋白质功能的突变是有害或中性的，这一点不用回避——把一台精密机器随机砸一锤，大概率砸坏而不是砸好。但“有害”永远是相对某个环境说的：镰刀型突变在杂合状态下对疟疾地区有利；让成年人继续消化乳糖的突变，在牧业社会是宝贝，在没奶喝的社会毫无意义；细菌的青霉素抗性在医院里是超能力，在没有药的环境里可能只是累赘。脱离环境谈突变的好坏，就像脱离天气谈雨衣的好坏。

\textbf{第四，突变率本身也在被进化塑造。} 修复太松，错误成灾；修得太死，代价高昂而且一点变化余地都不留。生命实际选择的，是一个“足够低但不为零”的中间值——这本身就是“突变既是风险也是材料”的最好注脚。

\begin{wuqu}
“突变都是有害的，所以突变只是错误，和进化扯不上关系。”——前半句就不成立。首先，大量突变是沉默的或落在无关紧要的位置，中性收场，你身上那几十个新突变里的大多数就属于此类。其次，反例一抓一大把：杂合状态的镰刀型突变在疟疾地区保护携带者；让成年人保持乳糖酶活性的突变，撑起了好几个大洲的奶业文明；细菌的一点点序列改动，能让它在抗生素里活下来（这对细菌是利好，对我们是麻烦——“有利”从来是对携带者自己说的）。正确的说法是：\textbf{多数突变中性或有害，少数在特定环境中有利}——而这少数，正是进化取之不尽的原材料。
\end{wuqu}

\section{回到那棵苹果树}

现在兑现开场的承诺。那根枝条的秘密是：在这棵树还是个芽的时候，某个细胞的 DNA 发生了一次突变，恰好影响了控制果皮色素（花青素）合成或分布的基因。这个细胞分裂、长大，长成了整根枝条——枝条上的每个细胞都带着这同一个突变，于是每一颗果子都呈现深红色。这叫\textbf{芽变}，是体细胞突变的经典案例。

为什么它年年如此、嫁接了也如此？因为无性繁殖不经过“重新洗牌”，枝条的每个细胞都忠实地带着那处改动。反过来，如果你把这颗红苹果的种子种下去，长出的树大概率退回普通模样——种子的细胞来自生殖过程，而那处突变多半没有进入生殖细胞。事实上，我们今天吃的许多知名水果品种，正是果农几十年上百年里不断“捉住”芽变、再用嫁接固定下来的成果：富士苹果的一系列深红色芽变品种，就是这么来的。

那盆金边吊兰也一样：绿色组织和白色组织的分界，就是当年那次突变发生的细胞 lineage 的分界。至于偶尔冒出的全白叶子——那是一个“两个 copies 都出问题”或关键突变更彻底的细胞后代，它彻底造不出叶绿素，无法进行光合作用，离开绿色部分的供养就活不下去。同一个基因，在叶片上是“好看的花纹”，在整株上是“致命的错误”。错误、材料还是创新？答案第三次回到同一个地方：看什么位置，看什么环境。

\section{错误、材料，还是创新}

绕了一圈，标题的问题可以回答了：突变\emph{同时}是这三者。

它是错误——从单个分子的角度看，它就是复制和损伤留下的瑕疵，修复系统全年无休地与之作战，多数改变功能的突变确实把事情搞砸了。它是材料——从群体和时间的角度看，每一代生物都携带着海量新序列，为任何可能到来的环境变化预备了候选答案。它偶尔是创新——当环境恰好选中某一处改动，昨天的错别字就变成了今天的新功能：能消化奶的肠道、能抵抗疟疾的血细胞、能躲开药物的细菌。

但请注意，突变本身只管“生产材料”，从来不负责“评选”。细菌不知道自己需要抗性，苹果树不知道自己可以更红——变异随机地出现，环境沉默地筛选。那个“评选”的过程有一条专门的名字和一整章的篇幅：第 76 章，自然选择。而工业界早已学会了人工加速“生产材料”这一步：把种子送上太空或送进辐照室提高突变率，再从成千上万株后代里挑选偶尔出现的优良个体——太空椒、辐照小麦就是这么来的。诱变负责出题，选择负责判卷，分工从四十亿年前到今天，一直没变。

\begin{tiaozhan}
波动实验真正说服人的地方在统计里。假设“接触后诱导”成立，每个细菌独立地以小概率变身，那么几十个小管里的抗性菌数应该服从泊松分布——它有个铁律：方差等于均值。而“接触前随机突变”成立时，早期突变会增殖出庞大克隆，方差会远远大于均值。卢里亚和德尔布吕克数出来的正是后者：方差比均值大了几十甚至上百倍。同样的数学也出现在别处： radioactive 衰变的计数涨落、工厂质检的抽样波动，都是泊松分布的地盘。学会问“波动该有多大”，你就有了一把能戳破很多伪科学论断的尺子。这部分跳过不影响主线。
\end{tiaozhan}

\section*{练一练}

\begin{sikaoti}
\item 一段正常序列是 ATG GAG CTT（每三个字母一个密码子）。分别写出：一个沉默突变、一个错义突变、一个无义突变各可能长什么样？再在第 4 个字母后插入一个 A，写出移码之后密码子怎样重新分组。（可查阅第 69 章的密码子表。）
\item 抄写实验中，你的“突变率”约为每千字 3 个错。照这个比例，如果让你徒手抄写一部 30 亿字母的“基因组”，会出多少个错？再想想细胞实际只错几十个，校对和修复系统一共把错误率压低了多少倍？
\item 画一幅信息流图：从“某皮肤细胞受到紫外线照射”开始，画出两条路径——损伤被修复、损伤未被修复并在下次复制时固定为突变——一直画到“几十年后可能的结果”。在图上标出防晒霜在哪一个环节起作用。
\item 一位果农发现某棵桃树的一根枝条结果特别甜。他想把这个性状保留下来，应该嫁接这根枝条，还是收集这颗桃子的种子播种？用体细胞突变和生殖系突变的区别解释你的选择。
\item 有人说：“细菌是为了抵抗抗生素才突变的。”用波动实验的逻辑反驳他，并预测：如果把从未接触过某抗生素的细菌培养物涂到含药培养皿上，会看到什么？
\item 全白的吊兰叶子美则美矣，却活不长；半绿半白的叶子却代代相传。用“有利有害取决于环境与位置”的观点解释，为什么花叶性状能在花市里一直见到。
\end{sikaoti}
